\documentclass[11pt]{article}
\usepackage[margin=1in]{geometry}
\usepackage{amsmath,amssymb,amsthm}
\usepackage{graphicx}
\usepackage{float}
\usepackage[T1]{fontenc}
\usepackage{lmodern}
\usepackage{microtype}
\usepackage[hidelinks]{hyperref}

\newtheorem{theorem}{Theorem}
\newtheorem{corollary}[theorem]{Corollary}
\newtheorem{lemma}[theorem]{Lemma}
\newcommand{\norm}[1]{\lVert#1\rVert}
\newcommand{\abs}[1]{\lvert#1\rvert}
\newcommand{\Ncount}{N_\alpha}
\newcommand{\Schatten}{\mathcal S}

\title{A Complete Counting Characterization of\
Schatten Rhaly Matrices for Every Positive Exponent}
\author{Full solution packet for arXiv:2503.21930}
\date{22 July 2026}

\begin{document}
\maketitle

\begin{center}
\fbox{\parbox{0.91\textwidth}{\textbf{Status: full solution likely valid.}
For every $q>0$, membership of a Rhaly matrix $R_\alpha$ in
$\Schatten^q(\ell^2)$ is characterized by the $q$-moment of the length of
the source paper's greedy $(\varepsilon,L)$-sequence.  Explicit two-sided
quasinorm estimates are proved.  In particular, this answers the source
paper's missing-characterization question for the entire range
$0<q\leq1$.}}
\end{center}

\section{Source question}

Bellavita, Dellepiane, and Stylogiannis study the Rhaly matrix
\begin{equation}\label{eq:rhaly}
 (R_\alpha x)(k)=\alpha_k\sum_{j=0}^k x(j),
 \qquad k\in\mathbb N,
\end{equation}
for $\alpha=(\alpha_k)\in\ell^2$ \cite{source}.  Their Theorem 1.3 gives a
dyadic coefficient characterization of $R_\alpha\in\Schatten^q$ when
$1<q<\infty$.  Theorem 5.5 gives the corresponding dyadic condition only as
a necessary condition for arbitrary $q>0$.  In the conclusion, on source
PDF page 36, they state:
\begin{quote}
``Theorem 5.5 only provides a necessary condition for the $q$-Schatten
class membership, for $0<q\leq1$. A complete characterization is missing.''
\end{quote}

\begin{figure}[H]
 \centering
 \includegraphics[width=0.98\textwidth]{figures/open_problem_crop.png}
 \caption{The endpoint open question in arXiv:2503.21930v3, source PDF
 page 36.}
\end{figure}

\section{The coefficient-side counting invariant}

For a natural interval $I$ and $x$ supported in $I$, put
\[
 S_x(k)=\sum_{j=0}^k x(j),\qquad
 \mu_\alpha(I)=\sum_{k\in I}\abs{\alpha_k}^2.
\]
When $\mu_\alpha(I)>0$, let
\[
 (S_x)_I=\frac{1}{\mu_\alpha(I)}
 \sum_{k\in I}S_x(k)\abs{\alpha_k}^2
\]
and define
\begin{equation}\label{eq:L}
 L_\alpha(I)=\sqrt2\sup_{\substack{\operatorname{supp}x\subseteq I\\
                                  \norm{x}_2\leq1}}
 \left(\sum_{k\in I}\abs{S_x(k)-(S_x)_I}^2
                         \abs{\alpha_k}^2\right)^{1/2}.
\end{equation}
Set $L_\alpha(I)=0$ if $\mu_\alpha(I)=0$, and define
$L_\alpha([a,\infty))$ by increasing finite truncations.  Formula
\eqref{eq:L} is the equivalent form of the source paper's interval invariant
given by its Lemma 3.6.

For $\varepsilon>0$, start with $c_0=0$ and recursively set
\begin{equation}\label{eq:greedy}
 c_{r+1}=\inf\{t\in\mathbb N:t>c_r,
          \ L_\alpha([c_r,t-1])>\varepsilon\}.
\end{equation}
If $c_N<\infty$ and $c_{N+1}=\infty$, the sequence has length $N$;
if it never terminates, its length is $\infty$.  Denote this length by
$\Ncount(\varepsilon)$.  This is entirely determined by the coefficients
$\alpha$.  It is an extended integer-valued Borel function of
$\varepsilon$: for fixed $m$, the event $\Ncount(\varepsilon)\geq m$ is a
countable union of conditions involving finitely many comparisons
$L_\alpha([a,b])>\varepsilon$.

For $q>0$, define the endpoint counting functional
\begin{equation}\label{eq:Phi}
 \Phi_q(\alpha)=q\int_0^\infty
   \varepsilon^{q-1}\Ncount(\varepsilon)\,d\varepsilon
 \in[0,\infty].
\end{equation}
Also set
\begin{equation}\label{eq:sigma}
 \sigma_{-1}=\abs{\alpha_0},\qquad
 \sigma_k=\left(\sum_{j=2^k}^{2^{k+1}-1}
                    (j+1)\abs{\alpha_j}^2\right)^{1/2},\qquad
 M_\alpha=\sup_{k\geq-1}\sigma_k.
\end{equation}
The source paper's Theorem 2.4 says that $M_\alpha<\infty$ exactly when
$R_\alpha$ is bounded.

\section{Proof intuition}

The greedy length already measures the whole singular-value distribution,
not merely compactness.  The source paper proves two estimates.  If the
$\varepsilon$-greedy sequence has length $N$, then the rank of a natural
piecewise-constant approximation is at most $2N+1$, so
$a_{2N+2}(R_\alpha)\leq\varepsilon/\sqrt2$.  Conversely, $N$ disjoint greedy
intervals force $a_N(R_\alpha)\geq\varepsilon/\sqrt2$.

Thus the number of approximation values above a threshold is trapped between
the same greedy length, with only a factor two in the count and fixed
rescalings of the threshold.  Integrating this distributional sandwich via
the layer-cake identity gives the complete $\Schatten^q$ criterion for every
$q>0$.  The endpoint obstruction in the source proof disappears because we
do not dominate the greedy count by an additional harmonic sum over dyadic
levels; that domination is precisely where the restriction $q>1$ entered.

\section{Distributional sandwich}

Let $a_m=a_m(R_\alpha)$ be the approximation numbers, indexed so that
$a_1=\norm{R_\alpha}$, and put
\[
 n_\alpha(t)=\#\{m\geq1:a_m>t\}.
\]
Both sides are allowed to be infinite.

\begin{lemma}[source estimates in counting form]\label{lem:sandwich}
For every $\varepsilon>0$,
\begin{equation}\label{eq:sandwich}
 \Ncount(\varepsilon)
 \leq n_\alpha\!\left(\frac{\varepsilon}{2\sqrt2}\right),
 \qquad
 n_\alpha\!\left(\frac{\varepsilon}{\sqrt2}\right)
 \leq2\Ncount(\varepsilon)+1.
\end{equation}
\end{lemma}

\begin{proof}
Suppose first that $\Ncount(\varepsilon)=N<\infty$.  If $N\geq1$, source
Lemma 4.3 gives
\[
 a_N\geq\frac{\varepsilon}{\sqrt2}
       >\frac{\varepsilon}{2\sqrt2}.
\]
Hence at least $N$ approximation numbers exceed
$\varepsilon/(2\sqrt2)$.  The first assertion is trivial when $N=0$.
If the greedy sequence is infinite, source Lemma 4.3 applies with every
finite $N$, so the first counting function is infinite as well.

For the other direction, source Lemma 4.2 gives
\[
 a_{2N+2}\leq\frac{\varepsilon}{\sqrt2}
\]
when the length is $N<\infty$.  Therefore at most $2N+1$ approximation
numbers are strictly larger than $\varepsilon/\sqrt2$.  If the length is
infinite, the claimed inequality is automatic.
\end{proof}

\section{Full characterization}

\begin{theorem}[all positive Schatten exponents]\label{thm:main}
Let $q>0$ and suppose $R_\alpha$ is bounded on $\ell^2$.  With the convention
that all quantities may equal $+\infty$,
\begin{align}
 \sum_{m=1}^\infty a_m(R_\alpha)^q
 &\leq \norm{R_\alpha}^q+2^{1-q/2}\Phi_q(\alpha),
 \label{eq:upper}\tag{A}\\
 \norm{R_\alpha}^q+\Phi_q(\alpha)
 &\leq\bigl(1+(2\sqrt2)^q\bigr)
       \sum_{m=1}^\infty a_m(R_\alpha)^q.
 \label{eq:lower}\tag{B}
\end{align}
Consequently,
\begin{equation}\label{eq:criterion}
 R_\alpha\in\Schatten^q(\ell^2)
 \quad\Longleftrightarrow\quad
 \Phi_q(\alpha)<\infty.
\end{equation}
\end{theorem}

\begin{proof}
For every nonincreasing nonnegative sequence $(a_m)$ and every $q>0$,
Tonelli's theorem gives the extended layer-cake identity
\begin{equation}\label{eq:layercake}
 \sum_{m=1}^\infty a_m^q
 =q\int_0^\infty t^{q-1}\#\{m:a_m>t\}\,dt.
\end{equation}

Integrating the first inequality in \eqref{eq:sandwich}, substituting
$t=\varepsilon/(2\sqrt2)$, and using \eqref{eq:layercake}, we obtain
\[
 \Phi_q(\alpha)
 \leq(2\sqrt2)^q\sum_{m=1}^\infty a_m^q.
\]
Since $\norm{R_\alpha}=a_1$, this proves \eqref{eq:lower}.

For the reverse estimate, put $\varepsilon=\sqrt2t$ in the second inequality
of \eqref{eq:sandwich}.  Since $n_\alpha(t)=0$ for
$t\geq\norm{R_\alpha}$, another application of \eqref{eq:layercake} gives
\begin{align*}
 \sum_{m=1}^\infty a_m^q
 &\leq 2q\int_0^{\norm{R_\alpha}}
       t^{q-1}\Ncount(\sqrt2t)\,dt
       +q\int_0^{\norm{R_\alpha}}t^{q-1}\,dt\\
 &\leq2^{1-q/2}\Phi_q(\alpha)+\norm{R_\alpha}^q.
\end{align*}
This is \eqref{eq:upper}, and the equivalence follows.
\end{proof}

Combining the theorem with the source boundedness criterion yields a statement
containing only coefficient-side quantities.

\begin{corollary}[complete coefficient characterization]\label{cor:complete}
Let $\alpha\in\ell^2$ and $q>0$.  Then
\[
 R_\alpha\in\Schatten^q(\ell^2)
 \quad\Longleftrightarrow\quad
 M_\alpha<\infty\ \text{ and }\ \Phi_q(\alpha)<\infty,
\]
where $M_\alpha$ and $\Phi_q(\alpha)$ are given explicitly by
\eqref{eq:L}--\eqref{eq:Phi} and \eqref{eq:sigma}.  In particular, this is a
complete characterization for every $0<q\leq1$.
\end{corollary}

\begin{proof}
Theorem 2.4 of the source gives $M_\alpha<\infty$ if and only if
$R_\alpha$ is bounded.  On that boundedness domain, Theorem
\ref{thm:main} applies.  Every $\Schatten^q$ operator is bounded, so the two
conditions are also necessary.
\end{proof}

\section{Verification, scope, and novelty}

The proof has no computational dependency.  The important endpoint details
are the strict threshold in $n_\alpha(t)$, the harmless factor two in the
first half of \eqref{eq:sandwich}, and the separate operator-norm term in
\eqref{eq:upper}.  That last term cannot simply be discarded: if only
$\alpha_0$ is nonzero, all greedy lengths vanish while $R_\alpha$ is a
nonzero rank-one operator.

The result gives an exact greedy-counting characterization, not a reduction
to the simpler dyadic Besov expression used by the source for $q>1$.  Thus it
answers the stated request for a complete characterization, while leaving
open whether a comparably simple dyadic coefficient formula exists at the
quasi-Banach endpoints.

A bounded novelty search on 22 July 2026 covered the run's four lightweight
indexes; the exact arXiv id and title; searches for Rhaly/terraced matrices
with Schatten endpoints, approximation numbers, and $(\varepsilon,L)$
sequences; the current source version; and the June 2026 follow-up
\href{https://arxiv.org/abs/2606.15875}{arXiv:2606.15875}.  That follow-up
extends the method to factorable matrices but still states its Schatten
characterizations only for $1<p<\infty$.  No later source stating the
counting criterion above or resolving the $0<q\leq1$ question was found.
Novelty confidence is moderate.

\textbf{Human-review recommendation.}  Audit the interpretation of the
source's length convention at $N=0$ and $N=\infty$, then verify that Lemmas
4.2 and 4.3 apply with precisely the approximation-number indexing used in
\eqref{eq:sandwich}.  Once that interface is confirmed, the remaining proof
is Tonelli integration.

\begin{thebibliography}{9}
\bibitem{source}
C. Bellavita, E. Dellepiane, and G. Stylogiannis,
\emph{Boundedness, compactness and Schatten class for Rhaly matrices},
arXiv:2503.21930v3 (2025), especially Definition 4.1, Lemmas 4.2--4.3,
Theorem 2.4, Theorem 5.5, and Section 7.
\end{thebibliography}

\end{document}
